\section{Tarefa 5.1: Montagem de um circuito combinacional com portas NAND2}


A Porta XOR foi construída seguindo o esquemático realizado pelo professor, ao término da montagem, foi testado todas as 4 combinações de entradas possíveis, chegando a seguinte tabela verdade

\begin{table}[H]
    \centering
    \caption{Tabela-verdade obtida para a saída da porta lógica XOR construída}
    \begin{tabular}{ccc}
        \toprule
        $A$ & $B$ & $F$ \\
        \midrule
        0   & 0   & 0   \\
        0   & 1   & 1   \\
        1   & 0   & 1   \\
        1   & 1   & 0   \\
        \bottomrule
    \end{tabular}
    \label{tab:tableXor}
\end{table}

A partir da tabela verdade construída, vemos que a saída F bate com o resultado esperado para uma porta XOR de duas entradas, dessa maneira:

\begin{equation}
    F = A \oplus B
\end{equation}